\documentclass[10pt,a4paper]{amsart}
\usepackage{amsmath,amssymb,amsthm,graphicx}
\usepackage{verbatim}
\newenvironment{reference}{\comment}{\endcomment}
%\newenvironment{reference}{}{}
\newenvironment{slogan}{\comment}{\endcomment}
\newenvironment{history}{\comment}{\endcomment}
% For commutative diagrams we use Xy-pic
\usepackage[all]{xy}

% We use 2cell for 2-commutative diagrams.
\xyoption{2cell}
% Theorem environments.
%
\theoremstyle{plain}
\newtheorem{theorem}[subsection]{Theorem}
\newtheorem{proposition}[subsection]{Proposition}
\newtheorem{lemma}[subsection]{Lemma}

\theoremstyle{definition}
\newtheorem{definition}[subsection]{Definition}
\newtheorem{example}[subsection]{Example}
\newtheorem{exercise}[subsection]{Exercise}
\newtheorem{situation}[subsection]{Situation}

\theoremstyle{remark}
\newtheorem{remark}[subsection]{Remark}
\newtheorem{remarks}[subsection]{Remarks}

\numberwithin{equation}{subsection}

% Macros
%
\def\lim{\mathop{\mathrm{lim}}\nolimits}
\def\colim{\mathop{\mathrm{colim}}\nolimits}
\def\Spec{\mathop{\mathrm{Spec}}}
\def\Hom{\mathop{\mathrm{Hom}}\nolimits}
\def\Ext{\mathop{\mathrm{Ext}}\nolimits}
\def\SheafHom{\mathop{\mathcal{H}\!\mathit{om}}\nolimits}
\def\SheafExt{\mathop{\mathcal{E}\!\mathit{xt}}\nolimits}
\def\Sch{\mathit{Sch}}
\def\Mor{\mathop{\mathrm{Mor}}\nolimits}
\def\Ob{\mathop{\mathrm{Ob}}\nolimits}
\def\Sh{\mathop{\mathit{Sh}}\nolimits}
\def\NL{\mathop{N\!L}\nolimits}
\def\CH{\mathop{\mathrm{CH}}\nolimits}
\def\proetale{{pro\text{-}\acute{e}tale}}
\def\etale{{\acute{e}tale}}
\def\QCoh{\mathit{QCoh}}
\def\Ker{\mathop{\mathrm{Ker}}}
\def\Im{\mathop{\mathrm{Im}}}
\def\Coker{\mathop{\mathrm{Coker}}}
\def\Coim{\mathop{\mathrm{Coim}}}

% Boxtimes
%
\DeclareMathSymbol{\boxtimes}{\mathbin}{AMSa}{"02}

%
% Macros for moduli stacks/spaces
%
\def\QCohstack{\mathcal{QC}\!\mathit{oh}}
\def\Cohstack{\mathcal{C}\!\mathit{oh}}
\def\Spacesstack{\mathcal{S}\!\mathit{paces}}
\def\Quotfunctor{\mathrm{Quot}}
\def\Hilbfunctor{\mathrm{Hilb}}
\def\Curvesstack{\mathcal{C}\!\mathit{urves}}
\def\Polarizedstack{\mathcal{P}\!\mathit{olarized}}
\def\Complexesstack{\mathcal{C}\!\mathit{omplexes}}
% \Pic is the operator that assigns to X its picard group, usage \Pic(X)
% \Picardstack_{X/B} denotes the Picard stack of X over B
% \Picardfunctor_{X/B} denotes the Picard functor of X over B
\def\Pic{\mathop{\mathrm{Pic}}\nolimits}
\def\Picardstack{\mathcal{P}\!\mathit{ic}}
\def\Picardfunctor{\mathrm{Pic}}
\def\Deformationcategory{\mathcal{D}\!\mathit{ef}}
\usepackage{cleveref}
\crefname{originallemma}{original lemma}{original lemmas}
\Crefname{originallemma}{Original lemma}{Original lemmas}
\makeatletter
\newlabel{lemma-cover-upstairs}{{lemma-cover-upstairs}{source}{}{sourcepointer-1}{}}
\newlabel{lemma-cover-upstairs@cref}{{[originallemma][0][]lemma-cover-upstairs}{[1][0][]source}{}{}{}}
\makeatother
\crefname{originallemma}{original lemma}{original lemmas}
\Crefname{originallemma}{Original lemma}{Original lemmas}
\makeatletter
\newlabel{lemma-existence-plus-flat-base-change}{{lemma-existence-plus-flat-base-change}{source}{}{sourcepointer-2}{}}
\newlabel{lemma-existence-plus-flat-base-change@cref}{{[originallemma][0][]lemma-existence-plus-flat-base-change}{[1][0][]source}{}{}{}}
\makeatother
\crefname{originalreference}{original reference}{original references}
\Crefname{originalreference}{Original reference}{Original references}
\makeatletter
\newlabel{spaces-morphisms-lemma-quasi-compact-scheme-theoretic-image}{{spaces-morphisms-lemma-quasi-compact-scheme-theoretic-image}{source}{}{sourcepointer-3}{}}
\newlabel{spaces-morphisms-lemma-quasi-compact-scheme-theoretic-image@cref}{{[originalreference][0][]spaces-morphisms-lemma-quasi-compact-scheme-theoretic-image}{[1][0][]source}{}{}{}}
\makeatother
\crefname{originalreference}{original reference}{original references}
\Crefname{originalreference}{Original reference}{Original references}
\makeatletter
\newlabel{stacks-properties-lemma-points-cartesian}{{stacks-properties-lemma-points-cartesian}{source}{}{sourcepointer-4}{}}
\newlabel{stacks-properties-lemma-points-cartesian@cref}{{[originalreference][0][]stacks-properties-lemma-points-cartesian}{[1][0][]source}{}{}{}}
\makeatother
\title{A point in the closure of the image of a morphism lies in the space of its scheme theoretic image}
\author{The Stacks Project Authors}
\date{stacks-project:stacks-morphisms.tex:local-snapshot-20261012; GNU FDL 1.2 or later}
\begin{document}
\maketitle
\noindent\textit{Editorial note.} All delivered mathematical text is the author's own, copied from the pinned Stacks Project chapter source sources/stacks-algebraic-stacks-morphisms/pinned-original-source.tex; no statement or proof step was rewritten, reordered or removed.\par\medskip
\noindent\textit{Editorial note.} Preamble provenance: the Stacks Project keeps one single shared preamble for every chapter; this card takes it directly from the staged sources/stacks-algebraic-stacks-morphisms/preamble.tex as three verbatim spans (lines 8-12, 14-18 and 154-217), so no derived file was written and nothing was added or redefined. The wrapper class and the base packages are supplied by the assembler.\par\medskip
\noindent\textit{Editorial note.} Bibliography provenance: neither the counted statement nor its proof contains a citation (this lane rescans the whole span before writing the card), so the card carries no bibliography and no entry was reconstructed.\par\medskip
\noindent\textit{Editorial note.} Selection provenance: this card was opened under the file-based control contract (CONTROL.md), with the mechanical pre-card gate run on its statement and proof interval before the card was written.\par\medskip
\noindent\textit{Allowed standard prerequisites (editorial).}\begin{itemize}
\item The behaviour of the scheme theoretic image under flat base change. Invoked inside the counted proof (source lines 9001-9004). Identifies the scheme theoretic image of the base-changed morphism with the base change of the image.
\item The reduction of the scheme theoretic image along a smooth cover of the source. Invoked inside the counted proof (source lines 9007-9009). Replaces the source of the base-changed morphism by a scheme.
\item The statement that the image of a quasi-compact scheme onto its scheme theoretic image is dense. Invoked inside the counted proof (source lines 9010-9012) as a cross-chapter pointer. Provides density of the image of the presenting scheme.
\item The description of the points of a fibre product of algebraic stacks. Invoked inside the counted proof (source lines 9005-9007) as a cross-chapter pointer. Identifies the inverse image of the space of the image with the space of the base-changed image.
\end{itemize}\medskip
\ifdefined\hypertarget\hypertarget{sourcepointer-1}{}\fi
\noindent\textit{Source reference.} Original lemma lemma-cover-upstairs: sources/stacks-algebraic-stacks-morphisms/pinned-original-source.tex, source lines 8754--8754. Navigation pointer only: the counted text refers to the Stacks item lemma-cover-upstairs, which is an existing labeled result of this same chapter and is not part of this card. The pointer carries the target's own label as its number and its definition line as its span; its content is not reproduced here.\par\medskip
\ifdefined\hypertarget\hypertarget{sourcepointer-2}{}\fi
\noindent\textit{Source reference.} Original lemma lemma-existence-plus-flat-base-change: sources/stacks-algebraic-stacks-morphisms/pinned-original-source.tex, source lines 8885--8885. Navigation pointer only: the counted text refers to the Stacks item lemma-existence-plus-flat-base-change, which is an existing labeled result of this same chapter and is not part of this card. The pointer carries the target's own label as its number and its definition line as its span; its content is not reproduced here.\par\medskip
\ifdefined\hypertarget\hypertarget{sourcepointer-3}{}\fi
\noindent\textit{Source reference.} Original reference spaces-morphisms-lemma-quasi-compact-scheme-theoretic-image: sources/stacks-algebraic-stacks-morphisms/pinned-original-source.tex, source lines 9010--9010. Navigation pointer only: the counted text refers to the Stacks item spaces-morphisms-lemma-quasi-compact-scheme-theoretic-image, which belongs to another chapter of the project and is not part of this card. The pointer carries the target's own label as its number because this chapter's staged snapshot does not contain the other chapter's tag table; its content is not reproduced here.\par\medskip
\ifdefined\hypertarget\hypertarget{sourcepointer-4}{}\fi
\noindent\textit{Source reference.} Original reference stacks-properties-lemma-points-cartesian: sources/stacks-algebraic-stacks-morphisms/pinned-original-source.tex, source lines 2126--2126. Navigation pointer only: the counted text refers to the Stacks item stacks-properties-lemma-points-cartesian, which belongs to another chapter of the project and is not part of this card. The pointer carries the target's own label as its number because this chapter's staged snapshot does not contain the other chapter's tag table; its content is not reproduced here.\par\medskip
\begin{lemma}
\label{lemma-topology-scheme-theoretic-image}
Let $f : \mathcal{X} \to \mathcal{Y}$ be a quasi-compact
morphism of algebraic stacks. Let $\mathcal{Z} \subset \mathcal{Y}$
be the scheme theoretic image of $f$. Then $|\mathcal{Z}|$
is the closure of the image of $|f|$.
\end{lemma}

\begin{proof}
Let $z \in |\mathcal{Z}|$ be a point.
Choose an affine scheme $V$, a point $v \in V$, and a smooth morphism
$V \to \mathcal{Y}$ mapping $v$ to $z$.
Then $\mathcal{X} \times_\mathcal{Y} V$ is a quasi-compact algebraic stack.
Hence we can find an affine scheme $W$ and a surjective smooth
morphism $W \to \mathcal{X} \times_\mathcal{Y} V$.
By Lemma \ref{lemma-existence-plus-flat-base-change}
the scheme theoretic image of
$\mathcal{X} \times_\mathcal{Y} V \to V$ is
$Z = \mathcal{Z} \times_\mathcal{Y} V$.
Hence the inverse image of $|\mathcal{Z}|$ in $|V|$ is $|Z|$ by
Properties of Stacks, Lemma \ref{stacks-properties-lemma-points-cartesian}.
By Lemma \ref{lemma-cover-upstairs} $Z$ is
the scheme theoretic image of $W \to V$.
By Morphisms of Spaces, Lemma
\ref{spaces-morphisms-lemma-quasi-compact-scheme-theoretic-image}
we see that the image of $|W| \to |Z|$ is dense.
Hence the image of $|\mathcal{X} \times_\mathcal{Y} V| \to |Z|$
is dense. Observe that $v \in Z$.
Since $|V| \to |\mathcal{Y}|$ is open, a topology argument
tells us that $z$ is in the closure of the image of $|f|$ as desired.
\end{proof}

\end{document}
